% Gerado por regressao/06_tabelas_latex.R (projeto Modelagem). Não editar à mão.
\begin{landscape}
\begin{table}[htbp]
\caption{Robustez R8 — Assistência e Previdência somadas (Sakurai, 2009)}\label{tab:robustez_R8}
\centering
\footnotesize
\setlength{\tabcolsep}{3pt}
\begin{tabular}{@{}>{\raggedright\arraybackslash}p{0.25\linewidth}>{\centering\arraybackslash}p{72pt}@{}}
\hline
\textbf{Variável} & \textbf{\hspace{0pt}Assistência e Previdência} \\ \hline
\multicolumn{2}{@{}l}{\textit{Modelo A (ciclo eleitoral)}} \\
Ano eleitoral & 10,004 \\
 & (6,913) \\
Ano pré-eleitoral & 10,708 \\
 & (6,655) \\
\hline
Observações & 70.515 \\
Municípios & 5.568 \\
R\textsuperscript{2} & 0,383 \\[6pt]
\multicolumn{2}{@{}l}{\textit{Modelo B (coalizões)}} \\
Prefeito na coalizão & 5,023\textsuperscript{***} \\
do governador & (1,627) \\
Prefeito na coalizão & $-$2,323 \\
do presidente & (2,339) \\
\hline
Observações & 70.515 \\
Municípios & 5.568 \\
R\textsuperscript{2} & 0,089 \\
\hline
\end{tabular}
\fonte{Elaboração própria.}
\nota{Modelo A: efeito fixo de município, erros-padrão de Driscoll-Kraay. Modelo B: efeitos fixos de município e de ano, erros-padrão agrupados por município. Erros-padrão entre parênteses. \textsuperscript{***}p$<$0,01; \textsuperscript{**}p$<$0,05; \textsuperscript{*}p$<$0,1.}
\end{table}
\end{landscape}
